Des biologistes étudient une population de bactéries. On modélise l'évolution de cette population par la fonction $f$ définie sur l'intervalle $[0 ; +\infty[$ par $f(t) = 10 \times 1,15^t$ où $t$ est exprimé en heure et $f(t)$ en \textbf{milliers} de bactéries.

L'expérience débute quand $t=0$.

\begin{enumerate}
	\item
	\begin{enumerate}
		\item Calculer le nombre de bactéries, exprimé en milliers, au début de l'expérience.
		\item Calculer le nombre de bactéries, exprimé en milliers, au bout d'une heure.
	\end{enumerate}
	\item Déterminer le sens de variation de la fonction $f$ sur $[0 ; +\infty[$ en justifiant.
	\item On a $f(1,5) \approx \np{12,332}$. Interpréter ce résultat dans le contexte de l'expérience.
\end{enumerate}
